\chapter{Una Cuadrícula Infinita en el Espacio: ¿Qué es un Retículo?}
\label{cap:reticulos_geometria}

\section{De la hoja cuadriculada al infinito}

Toma un cuaderno cuadriculado escolar. Cada intersección de dos líneas es un punto $(x, y)$ donde las coordenadas $x$ e $y$ son números enteros:
\[
(0,0),\; (1,0),\; (0,1),\; (2,3),\; (-1, 4),\dots
\]
Ese conjunto de puntos discretos y perfectamente espaciados que se extiende hacia el infinito se llama en matemáticas $\mathbb{Z}^2$.

Ahora imagina que deformamos esa cuadrícula: en lugar de caminar en pasos perpendiculares de longitud $1$ hacia la derecha y hacia arriba, elegimos dos flechas (vectores) cualesquiera en el plano, digamos:
\[
\mathbf{b}_1 = \begin{pmatrix} 2 \\ 1 \end{pmatrix}, \quad \mathbf{b}_2 = \begin{pmatrix} 1 \\ 3 \end{pmatrix}
\]

\begin{figure}[htbp]
\centering
\begin{tikzpicture}[scale=0.85]
  % Puntos del reticulo
  \foreach \x in {-2,...,3} {
    \foreach \y in {-2,...,2} {
      \coordinate (P) at ({\x*1.5 + \y*0.7}, {\x*0.4 + \y*1.8});
      \fill[PQCMainBlue!80!black] (P) circle (2.2pt);
    }
  }
  
  % Origen
  \coordinate (O) at (0,0);
  \node[below left, font=\small\bfseries] at (O) {$(0,0)$};
  
  % Vectores base b1 y b2
  \draw[-{Stealth[length=3mm]}, line width=1.8pt, PQCTeal] (O) -- (1.5, 0.4) node[midway, below right, font=\small\bfseries] {$\mathbf{b}_1$};
  \draw[-{Stealth[length=3mm]}, line width=1.8pt, PQCGreen] (O) -- (0.7, 1.8) node[midway, above left, font=\small\bfseries] {$\mathbf{b}_2$};
  
  % Combinacion de ejemplo
  \draw[-{Stealth[length=2.5mm]}, line width=1.2pt, PQCPurple, dashed] (O) -- ({2*1.5 + 1*0.7}, {2*0.4 + 1*1.8}) node[above right, font=\small\bfseries] {$2\mathbf{b}_1 + \mathbf{b}_2$};

\end{tikzpicture}
\caption{Un retículo en 2 dimensiones generado por los vectores $\mathbf{b}_1$ y $\mathbf{b}_2$. Cada punto es el resultado de sumar múltiplos enteros de las flechas base.}
\label{fig:reticulo_2d}
\end{figure}

Si sumamos o restamos estas dos flechas multiplicadas por números enteros (por ejemplo $3\mathbf{b}_1 - 2\mathbf{b}_2$), generamos una red infinita de puntos regulares en el plano.

\begin{intuicion}[Definición para primer semestre: ¿Qué es un Retículo?]
Un \textbf{retículo} (\textit{lattice} en inglés) es una cuadrícula infinita de puntos en el espacio donde:
\begin{enumerate}
  \item El origen $(0,0,\dots,0)$ siempre forma parte de la cuadrícula.
  \item Si tomas dos puntos cualquiera de la red y los sumas, caes exactamente en otro punto de la red.
  \item Si caminas de un punto al siguiente, siempre hay una distancia mínima positiva (los puntos están separados, no se amontonan de manera continua como el agua).
\end{enumerate}
\end{intuicion}

\section{Bases buenas y bases malas: El poder de la perspectiva}

Aquí viene el primer gran secreto de los retículos. Un mismo retículo (la misma constelación idéntica de puntos) puede ser generado por \textbf{diferentes pares de vectores base}.

\begin{itemize}
  \item \textbf{Una base ``buena'':} Está formada por vectores cortos y casi perpendiculares (a casi 90 grados). Es como tener una regla bien calibrada: es fácil saber en qué celda estás y a qué punto estás más cerca.
  \item \textbf{Una base ``mala'':} Está formada por vectores larguísimos, muy inclinados y casi paralelos entre sí. Aunque generan la misma red, visualmente parecen un laberinto enmarañado.
\end{itemize}

\begin{figure}[htbp]
\centering
\begin{tikzpicture}[scale=0.85]
  % Puntos
  \foreach \x in {-2,...,4} {
    \foreach \y in {-2,...,3} {
      \coordinate (P) at ({\x*1.3 + \y*0.3}, {\x*0.2 + \y*1.4});
      \fill[gray!60] (P) circle (1.8pt);
    }
  }
  
  \coordinate (O) at (0,0);
  
  % Base buena (verde)
  \draw[-{Stealth[length=2.5mm]}, line width=1.6pt, PQCGreen] (O) -- (1.3, 0.2) node[below, font=\footnotesize\bfseries] {$\mathbf{b}_1$ (Corta)};
  \draw[-{Stealth[length=2.5mm]}, line width=1.6pt, PQCGreen] (O) -- (0.3, 1.4) node[left, font=\footnotesize\bfseries] {$\mathbf{b}_2$ (Corta)};
  
  % Base mala (rojo)
  \draw[-{Stealth[length=2.5mm]}, line width=1.6pt, PQCCrimson] (O) -- ({3*1.3 + 2*0.3}, {3*0.2 + 2*1.4}) node[above right, font=\footnotesize\bfseries] {$\mathbf{v}_1$ (Larga y distorsionada)};
  \draw[-{Stealth[length=2.5mm]}, line width=1.6pt, PQCCrimson] (O) -- ({2*1.3 + 1*0.3}, {2*0.2 + 1*1.4}) node[midway, below right, font=\footnotesize\bfseries] {$\mathbf{v}_2$};

\end{tikzpicture}
\caption{La misma cuadrícula vista a través de dos bases distintas: la base buena (verde) es clara y ordenada; la base mala (roja) es confusa y distorsionada.}
\label{fig:base_buena_mala}
\end{figure}

\begin{metafora}[La cuadrícula de una ciudad]
Imagina una ciudad como Manhattan, con avenidas y calles perpendiculares (base buena). Si te dicen ``avanza 3 cuadras al norte y 2 al este'', sabes exactamente dónde estás y cuál es la casa más cercana. 

Ahora imagina que te dan dos direcciones raras: una avenida diagonal que avanza 100 kilómetros hacia el noreste y otra que avanza 99 kilómetros casi en la misma dirección (base mala). Para llegar a la tienda de la esquina tendrías que sumar y restar números gigantescos con precisión astronómica. Cualquier pequeño error te dejaría en otra provincia.
\end{metafora}

\section{El problema geométrico central: CVP (El vector más cercano)}

Ahora presentemos el problema que le quita el sueño a las supercomputadoras: el \textbf{Problema del Vector Más Cercano} (\textit{Closest Vector Problem} o CVP).

\begin{intuicion}[El juego del CVP]
\textbf{Planteamiento:} Imagina que te entregan un punto en el espacio $\mathbf{t}$ que no pertenece al retículo (cayó en el ``espacio vacío'' entre los puntos de la cuadrícula, por culpa de un pequeño error). 

\textbf{El desafío:} ¿Cuál es el punto del retículo más cercano a $\mathbf{t}$?
\end{intuicion}

\begin{figure}[htbp]
\centering
\begin{tikzpicture}[scale=0.95]
  % Puntos
  \foreach \x in {-1,...,3} {
    \foreach \y in {-1,...,2} {
      \coordinate (P) at ({\x*1.8 + \y*0.4}, {\x*0.3 + \y*1.6});
      \fill[PQCMainBlue] (P) circle (2.2pt);
    }
  }
  
  % Punto objetivo t (fuera de la malla)
  \coordinate (T) at ({1*1.8 + 1*0.4 + 0.35}, {1*0.3 + 1*1.6 + 0.45});
  \fill[PQCCrimson] (T) circle (3pt) node[above right, font=\small\bfseries] {Punto objetivo $\mathbf{t}$ (con ruido)};
  
  % Punto mas cercano
  \coordinate (Closest) at ({1*1.8 + 1*0.4}, {1*0.3 + 1*1.6});
  \draw[PQCGreen, line width=1.5pt] (Closest) circle (6pt);
  \draw[-{Stealth[length=2.5mm]}, line width=1.2pt, PQCAmber] (Closest) -- (T) node[midway, below right, font=\footnotesize\bfseries] {ruido $\mathbf{e}$};
  \node[below left, font=\footnotesize\bfseries, PQCGreen] at (Closest) {¡Punto más cercano!};
\end{tikzpicture}
\caption{El problema CVP: Dado un punto $\mathbf{t}$ que fue desplazado por un ruido $\mathbf{e}$, el objetivo es encontrar cuál punto azul era el original.}
\label{fig:cvp}
\end{figure}

\begin{itemize}
  \item \textbf{Si tienes la base buena:} Miras tus coordenadas, redondeas al entero más cercano, y en un instante descubres el punto exacto.
  \item \textbf{Si solo tienes la base mala:} No sabes hacia dónde redondear. En dimensiones altas (por ejemplo $n = 500$ o $n = 768$ dimensiones), hay tantas posibilidades de búsqueda que ni siquiera todas las computadoras del planeta trabajando juntas durante siglos podrían encontrar el punto correcto.
\end{itemize}

¡Eureka! Ya tenemos nuestra primera trampa criptográfica:
\begin{itemize}
  \item \textbf{Tu Clave Privada} es la \textbf{base buena} (corta y ordenada).
  \item \textbf{Tu Clave Pública} es la \textbf{base mala} (larga y desordenada).
  \item Cualquiera puede usar tu clave pública para sumar vectores y agregarle un poco de ruido (cifrar).
  \item Solo tú, con tu base buena, puedes quitar el ruido y encontrar el mensaje original (descifrar).
\end{itemize}

\begin{resumencapitulo}
\begin{itemize}
  \item Un retículo es una cuadrícula infinita y periódica de puntos generada por combinaciones lineales enteras de vectores.
  \item Un mismo retículo tiene infinitas bases: bases ``buenas'' (vectores cortos, fáciles de manejar) y bases ``malas'' (vectores largos y casi paralelos).
  \item El Problema del Vector Más Cercano (CVP) consiste en hallar el punto del retículo más próximo a un punto ruidoso.
  \item En cientos de dimensiones, CVP es extremadamente fácil si conoces la base buena, pero computacionalmente imposible si solo tienes la base mala.
\end{itemize}
\end{resumencapitulo}
